\section{Implications for reductions and compression}\label{sec:discussion}
The moment formulas connect the input distribution, block partition, and
linear operator to the output error. For identical coordinate laws, the
weight sum determines how strongly the common error component enters a
reduction. Balanced weights remove the bias and the term proportional to
$(\mathbf1^Tw)^2$, leaving $(v-\kappa)\|w\|_2^2$. Positive weights can
instead accumulate both covariance and bias. This explains why stochastic
rounding can improve a positive-weight reduction while increasing the
mean-square error of a balanced one: its larger marginal variance is then
the deciding term. The appropriate rounding rule depends on the weights
and the accuracy criterion.

Partitioning changes both the selected scales and the covariance structure.
Each candidate partition therefore requires its own moments. The storage
cost of additional scales also changes the available coefficient precision.
In an elliptic solve, block layout further determines how the covariance
aligns with the Green operator. A model that retains only marginal error
variances loses this interaction, even when it estimates those variances
accurately. The equal-budget experiments demonstrate how the complete
moments predict solution error and guide the comparison of prescribed
compression choices. The calibration experiment shows how this workflow
degrades when the source law is estimated from a separate sample.

Block length and total problem size have different effects. If $N=Kn$
independent inputs occupy $K$ identical blocks of fixed length $n$, their
total error satisfies
\begin{equation}\label{eq:multiblock-mse}
 \mathbb E\left[\left(\sum_{i=1}^{N}e_i\right)^2\right]
 =Nv\mathcal I_n+N^2\mu^2.
\end{equation}
The centered standard deviation grows as $\sqrt N$, while a nonzero mean
error contributes quadratically to mean-square error. An inflation factor
describes the centered component; accuracy assessment requires both terms.

The exponential counterexample concerns growth of the block itself. At
fixed grid precision, its increasing maximum makes each fixed coordinate
eventually round to zero. The error approaches the negative input, and
pair correlation vanishes, but multiplication by $n-1$ still permits
diverging inflation. The dyadic location of the maximum produces pronounced
changes with block length. The exact finite-block evaluator resolves these
changes without extrapolating a power law from one phase.

Latent factors separate a structured field's common fluctuations from its
independent innovations. Conditional moment evaluation preserves dependence
between spatial blocks through the varying mean error. More general joint
laws require joint truncated moments; data-dependent weights require mixed moments of
the weights and quantization errors. Quantizing both operands and including
product or accumulation rounding adds further error terms. Saturation,
nonuniform element grids, and finite random bits can be incorporated by
supplying the corresponding scalar quantizer moments
\cite{elarar2024}.
The applications use a uniform integer-mantissa grid with an 8-bit signed
exponent and binary64 arithmetic. They establish accuracy comparisons at a
fixed payload budget for specified source laws and discrete operators.
Transfer to measured scientific fields, nonuniform MX element grids,
saturation, and low-precision accumulation requires separate validation.
Encoding cost, memory traffic, and solver speed were not measured.

\section{Conclusions}\label{sec:conclusions}
Maximum-selected block quantization creates two competing covariance terms:
variation of the conditional error mean and the correction caused by
excluding an all-small block. One-dimensional truncated moments compute
both exactly and give variance and mean-square error predictions for
weighted reductions. For identical input laws, dropping the conditional
pair correction has a distribution-free additive error bound of one in
the variance inflation factor. The heterogeneous and conditional-mixture
extensions supply the covariance for spatially varying and correlated sources.

The support limits and exponential counterexample show why fixed-pair
decorrelation does not imply unit inflation. The experiments verify the
finite-block identities and connect the resulting moments to weight
balance, storage allocation, spatial layout, and elliptic solution error.
Together, the analysis and evaluator turn a statistical quantization model
into a computable accuracy prediction for a specified input law and
compression design.
